\section{Analýza stávající aplikace}

Existující řešení je webová aplikace zaměřena na výuku anglického jazyka. Frontendová část aplikace je naprogramována pomocí technologií NextJS a React. Dále využívá knihovny, jako například Material UI. Celé uživatelské rozhraní je navrženo v nástroji Figma. %% write - přidat citaci na bakalářku

S backendovou částí frontend komunikuje přes REST API. Kromě aplikace existuje i kompletní mock backendového REST API, který byl vytvořen pomocí nástroje Mockoon. %% write - tady asi zacitovat bakalářku

V následujících částech textu se zaměřím nejprve na analýzu existujícího kódu a architektury a následně na jednotlivé funkcionality aplikace.

\subsection{Analýza kódu a architektury}

%% Momentální stav

\subsubsection{Design system}

Aplikace využívá knihovnu Material UI, %% write - přidat citaci
která poskytuje základní komponenty uživatelského rozhraní. Komponenty jsou přizpůsobeny především globálními styly tak, aby odpovídaly navrženému designu v nástroji Figma. Ze základních komponentů jsou následně sestaveny komplexnější komponenty, které jsou využívány v aplikaci. 

V kódu je také zavedena knihovna Storybook, která umožňuje snadné prohlížení, dokumentaci a testování jednotlivých komponentů.

Celková struktura komponentů je přehledná, nicméně stojí za úvahu zdokonalit tuto strukturu vytvořením samostatného systému komponentů, který by byl oddělen od aplikace a to hlavně z následujících důvodů.

Hlavním důvodem pro vytvoření samostatné knihovny komponentů je fakt, že růstem projektu je třeba zachovat konzistenci vzhledu značky a uživatelského rozhraní napříč různými částmi projektu, kterými jsou kromě této aplikace i například webové stránky a vzdělávací a marketingové materiály. Jelikož projekt v rámci marketingu využívá videa a materiály, která jsou také vytvářena pomocí React komponentů, začíná být centrální systém komponentů nutný.

Díky vytvoření samostatného deisgn systému by bylo možné izolovat základní komponenty a jejich styly od zbytku aplikace. To by umožnilo snadnější návrh, verzování a údržbu designu celé aplikace. Navíc by se zajistila konzistence vzhledu a funkčnosti komponentů napříč všemi částmi projektu.

\subsubsection{Kvalita kódu}

% Write - sem někam napsat, že v kódu nejsou žádné automatizované testy

Při analýze jsem nalezl několik příležitostí ke zlepšení celkové kvality kódu. V průběhu let používané knihovny a nástroje prošly aktualizacemi na nové verze, které nabízejí nové funkce a možností. Jejich aktualizace by mohla přispět ke zvýšení kvality kódu, bezpečnosti, výkonu a celkové udržitelnosti celé aplikace.

Společně s knihovnami by bylo vhodné aktualizovat konvence a pravidla pro knihovny ESLint a Prettier a provést refaktoring kódu tak, aby odpovídal novým pravidlům.

Nakonec by bylo vhodné zvážit dockerizaci projektu, která by mohla přispět ke zjednodušení vývoje a zajištění konzistence vývojových prostředí.

\subsubsection{Zavedení umělé inteligence pro vývoj}

V průběhu posledních let došlo k výraznému rozvoji nástrojů využívajících umělou inteligenci k vývoji softwaru. Jelikož aplikace vznikala v době, kdy tyto nástroje nebyly zdaleka tak rozšířené, nejsou v aplikaci vůbec zavedeny. Proto by bylo vhodné tyto nástroje zavést do projektu tak, aby se zefektivnil celkový vývoj aplikace.

% Write - přidat se subsubsection o tom, že v aplikaci je Mock API a že by mohlo být vhodnější napojit na skutečné testovací API a že by API mohlo být generování automaticky podle openapi specifikace API - formulovat to jako "příležitost ke zlepšení" a říct proč mi to přijde lepší než mock API

